\documentclass[12pt,a4paper]{ctexart}
\usepackage{amsmath,amssymb}
\usepackage{booktabs}
\usepackage{array}
\usepackage{geometry}
\usepackage{enumitem}
\usepackage[hidelinks]{hyperref}
\geometry{margin=2.4cm}
\setlength{\parskip}{0.4em}
\setlength{\parindent}{2em}
\linespread{1.25}
\renewcommand{\arraystretch}{1.2}
\setcounter{secnumdepth}{-1}
\ctexset{
  section = {format = {\heiti \zihao{3}}},
  subsection = {format = {\heiti \zihao{4}}},
}

\begin{document}

\begin{center}
{\heiti\zihao{-2} A 题分学科学习路线图\par}
\vspace{0.4em}
{\heiti\zihao{4} ——从高中理化生到大一数学，到读懂这篇论文\par}
\end{center}

\section{一、这张图怎么用}

\textbf{你的起点}：高中物理、化学、生物 + 大一数学（导数、泰勒、牛顿迭代；ODE/PDE 只知定义）。\textbf{目标}：读懂并手推《A题论文》与教程《A题方法详解》中的全部方法。

三条原则：

\begin{enumerate}[leftmargin=2em]
\item \textbf{物理直觉先行，数学只是语言}。每个模块都从你高中学过的现象出发（热传导、扩散、蒸发、饱和汽压、活化能、自由水/结合水），再补数学形式——不要先啃公式。
\item \textbf{只在用到的地方学}。传热学、化工原理是大三教材，\textbf{不系统学}，只读本图指定的章节页码；主线永远是教程《A题方法详解》（下称"教程"）。
\item \textbf{手推才算学会}。每个模块配教程附录 A 的 12 道手推练习（编号 A1--A12），公式必须自己推过一遍。
\end{enumerate}

\textbf{总路线}（学习顺序即教程章序，本图告诉你每一章属于哪个学科、缺什么高中背景、配哪本教材哪一节）：

$$\text{物理定律(直觉)}\to\text{数学工具(语言)}\to\text{PDE 建模}\to\text{数值方法}\to\text{解析解}\to\text{化学原理(参数)}\to\text{工程干燥(解释)}\to\text{论文}$$

\section{二、起点盘点}

\begin{center}
{\zihao{5}\begin{tabular}{p{2.6cm}p{5.2cm}p{5.6cm}}
\toprule
学科 & 你已有的（高中/大一） & 本图帮你补到 \\
\midrule
数学 & 导数、泰勒展开、牛顿迭代、数列递推；矩阵消元入门 & 偏导与链式法则、三对角矩阵、级数；PDE 建模与三类定解条件；差分/稳定性/追赶法等数值方法 \\
物理 & 热传递三方式（初中）、比热容、分子扩散现象（选修 3-3）、饱和汽压与相对湿度（3-3）、蒸发吸热 & 傅里叶导热定律、菲克扩散定律、对流边界（牛顿冷却公式）；湿球温度（热质同时传递）；$\sqrt{Dt}$ 数量级估计 \\
化学 & 溶液浓度与质量分数、化学平衡、活化能（定性，选必 1） & 干基/湿基含水率、Arrhenius 定量式（$E_a$ 手推）、饱和蒸气压与相平衡 \\
生物 & 自由水/结合水（必修 1） & 干燥学的完整水分分类（结合/非结合/平衡/自由/临界） \\
\bottomrule
\end{tabular}}
\end{center}

\section{三、你没学过的概念速查表（一嘴版）}

论文与教程中出现的术语，每条一句话讲清（标 $\blacktriangle$ 的是高中学过、只需唤醒的）；更详细的定义见教程附录 B（33 条）与本表指出的教材章节。水分分类五条（非结合水、结合水、平衡水、自由水、临界含水量）按化工原理下册 14.2--14.3 节的标准定义给出。

\subsection{3.1 数学与数值}

\begin{center}
{\zihao{5}\begin{tabular}{p{2.7cm}p{8.3cm}p{3.4cm}}
\toprule
概念 & 一句话定义 & 出处 \\
\midrule
偏导数 & 只对一个变量求导、其余变量当常数 & 教程 1.1 \\
链式法则 & 复合函数求导的传递规则（$\dfrac{dz}{dx}=\dfrac{dz}{du}\dfrac{du}{dx}$ 的多元版） & 教程 1.1 \\
中心差商 & 用左右两点对称近似导数，误差是二阶的 & 教程 1.2 \\
截断误差 & 用差分代替导数时被舍去的泰勒余项 & 教程 1.2、4.5 \\
三对角矩阵 & 只有主对角线及相邻两条对角线非零的矩阵 & 教程 1.5 \\
追赶法（Thomas） & 专门解三对角方程组的消元法，工作量 $O(N)$ & 教程 7 \\
显式格式 & 新时刻的值直接由旧时刻算出 & 教程 4.3 \\
隐式格式 & 新时刻的值需要联立解一次方程组 & 教程 4.3 \\
数值稳定性 & 误差不随推进被放大失控的性质 & 教程 4.4 \\
放大因子 & 一个频率分量经一步计算被放大/缩小的倍数 $\xi$ & 教程 4.4 \\
控制体 & 每个网格节点"管辖"的一小块区域 & 教程 6.2 \\
守恒型离散 & 每块区域严格"流入$-$流出$=$积累"的离散方式 & 教程 6 \\
不动点迭代 & 反复代入 $x=g(x)$ 逼近解；Picard 迭代是一例 & 教程 8.2 \\
随体坐标 $\zeta=r/R(t)$ & 跟随表面一起伸缩的坐标，消除收缩边界的麻烦 & 教程 9.2 \\
保形插值（Pchip） & 过数据点的光滑曲线，不产生虚假凸起 & 教程 9.6 \\
分离变量法 & 设 $u(r,t)=f(r)g(t)$ 代入，把 PDE 拆成两个常微分方程 & 教程 11.1 \\
特征方程 & 使边值问题有非零解时 $\beta$ 必须满足的方程 & 教程 11.1 \\
Bessel 函数 $J_0,J_1$ & 柱坐标分离变量自然出现的特殊函数 & 教程 11.1 \\
正交展开 & 像坐标轴一样用一族函数展开任意函数来求系数 & 教程 11.1 \\
洛必达法则 & 求 $0/0$ 型极限的法则 $\blacktriangle$ & 教程 3.3 \\
\bottomrule
\end{tabular}}
\end{center}

\subsection{3.2 物理：传热、扩散与湿空气}

\begin{center}
{\zihao{5}\begin{tabular}{p{2.7cm}p{8.3cm}p{3.4cm}}
\toprule
概念 & 一句话定义 & 出处 \\
\midrule
傅里叶定律 & 热流量正比于温度梯度，方向从高温到低温 & 教程 2.1 \\
菲克定律 & 扩散通量正比于浓度梯度，方向从高浓到低浓 & 教程 2.1 \\
导热系数 $k$ & 材料导热本领的材料参数 & 教程 2.1 \\
热扩散率 $\alpha=k/(\rho c_p)$ & 温度"传播"的快慢 & 教程 2.2 \\
扩散系数 $D$ & 浓度"传播"的快慢 & 教程 2.3 \\
对流换热系数 $h$ & 表面与气流换热强度的系数（牛顿冷却公式的比例常数） & 教程 3.2 \\
对流传质系数 $h_m$ & 表面与气流水分交换强度的系数 & 教程 3.2 \\
三类边界条件 & 第一类给温度、第二类给热流、第三类给对流平衡 & 教程 3.1 \\
Robin 条件 & "表面导热通量$=$对流带走通量"的平衡式（第三类） & 教程 3.2 \\
Biot 数 & 内部导热阻力与表面换热阻力之比，$Bi=hR/k$ & 论文 7.1 \\
$\sqrt{Dt}$ & 扩散在时间 $t$ 内大致到达的深度 & 教程 1.4 \\
$R^2/\alpha$ & 热传导贯穿整个圆柱的特征时间 & 教程 12.4 \\
饱和蒸汽压 $p_s(T)$ & 该温度下空气中水蒸气的分压上限 $\blacktriangle$ & 教程 10.4 \\
相对湿度 RH & 实际水蒸气分压 $\div$ 饱和蒸汽压 $\blacktriangle$ & 教程 10.4 \\
汽化潜热 & 1 kg 水蒸发吸收的热量（只相变、不升温）$\blacktriangle$ & 教程 10.5 \\
显热 & 使温度改变的热量 & 教程 10.5 \\
湿球温度 $t_w$ & 湿润表面因蒸发冷却达到的动态平衡温度，只取决于空气状态 & 教程 10.4 \\
绝热饱和温度 & 空气绝热增湿到饱和时的温度；空气--水系统约等于湿球温度 & 化工原理下册 13.2.2 \\
热质同时传递 & 传热与传质互相影响、同时进行的耦合过程 & 化工原理下册第 13 章 \\
边界层 & 壁面附近流速/温度/浓度剧烈变化的薄层 & 化工原理上册第 6 章 \\
\bottomrule
\end{tabular}}
\end{center}

\subsection{3.3 化学与化工（干燥）}

\begin{center}
{\zihao{5}\begin{tabular}{p{2.7cm}p{8.3cm}p{3.4cm}}
\toprule
概念 & 一句话定义 & 出处 \\
\midrule
干基含水率 $C$ & 水质量 $\div$ 绝干物料质量（论文统一用干基） & 教程 10.1 \\
湿基含水率 & 水质量 $\div$ 物料总质量 & 教程 10.1 \\
表观活化能 $E_a$ & Arrhenius 指数中的"能垒"参数，由数据拟合得出（活化能概念高中已学 $\blacktriangle$） & 教程 10.2 \\
Arrhenius 公式 & 速率常数 $\propto\exp(-E_a/RT)$：温度越高、越过能垒越快 & 教程 10.2 \\
吸附/解吸 & 水分子附着于固体表面 / 从表面脱离 & 化工原理下册 14.2 \\
吸湿等温线 & 平衡含水率随相对湿度变化的曲线 & 任迪峰 3.2 节 \\
干燥速率曲线 & 干燥速率随含水率变化的曲线，分预热/恒速/降速段 & 化工原理下册 14.3 \\
恒速段 & 表面被非结合水覆盖、速率不变、表面温度$=$湿球温度 & 化工原理下册 14.3 \\
降速段 & 速率不断下降（内部扩散控制）、物料温度逐渐升高 & 化工原理下册 14.3 \\
临界含水量 $X_c$ & 恒速段与降速段的分界含水率 & 化工原理下册 14.3 \\
非结合水分 & 机械结合（表面润湿水、大孔隙水），蒸汽压$=$纯水饱和值，最易除去 & 化工原理下册 14.2 \\
结合水分 & 物理化学结合（吸附水、毛细凝聚水、细胞壁内水），蒸汽压$<$饱和值，较难除去 & 化工原理下册 14.2 \\
平衡水分 $X^*$ & 指定空气条件下干燥能达到的最低含水量，干燥除不掉 & 化工原理下册 14.2 \\
自由水分 & 超过平衡水分、可被干燥除去的水分 $=X-X^*$ & 化工原理下册 14.2 \\
导湿温性 & 温度梯度也驱动水分迁移的效应（雷科夫理论）；本文模型不含 & 任迪峰 4.2 \\
发汗 & 干燥后期停止加热、靠内部水分外移使药材回软的工序 & 任迪峰 \\
拟定态 & 对缓慢变化过程每一时刻按稳态近似处理 & 化工原理上册 6.6.5 \\
\bottomrule
\end{tabular}}
\end{center}

\textbf{易混提醒}：① 干燥学的"自由水分"（按平衡定义）$\ne$ 生物学的"自由水"（按结合方式）；② 两套分类的关系：自由水分 $=$ 全部非结合水分 $+$ 部分结合水分，RH$=100\%$ 时平衡水分$=$结合水分；③ 生化"结合水"靠氢键/离子键与蛋白质结合，\textbf{配位键}对应的是结晶水（如 $\mathrm{CuSO_4\cdot5H_2O}$），属于结合水分的最极端情况，常规干燥除不掉。

\section{四、学科一：数学（三种身份——语言、建模、求解）}

\subsection{模块 1\quad 微积分工具箱（P0，必须先学）}

\textbf{高中锚点}：导数是"变化率"、泰勒是"用多项式逼近函数"（大一已学）。

\textbf{新学内容}：① 偏导数与多元链式法则（教程 1.1 节）；② 泰勒展开导出两个中心差商及截断误差系数 $1/6$、$1/12$（1.2 节）；③ 通量与守恒的积分思想（1.3 节）；④ 三对角矩阵的形态（1.5 节）与追赶法消元（7.2 节）。

\textbf{资料位置}：教程 1.1--1.5、7.1--7.4 节；同济《高等数学》多元函数微分学、无穷级数两章（只补偏导与幂级数）。

\textbf{手推练习}：A2（两个差商系数）、A8（追赶法公式 $+$ $4\times4$ 例题）。

\textbf{学完标准}：不看提示能写出中心差商公式并给出截断误差首项；能解释追赶法为什么是 $O(N)$ 而不是 $O(N^3)$。

\subsection{模块 2\quad PDE 建模与定解条件（P0，必须先学）}

\textbf{高中锚点}：守恒（进多少 $-$ 出多少 $=$ 攒多少）；初中热传导定性描述。

\textbf{新学内容}：① 由壳层守恒手推柱坐标热传导方程与水分扩散方程（教程 2.2、2.3 节，含 $1/r$ 的来历）；② 三类边界条件与 Robin 条件的来历（3.1、3.2 节）；③ $r=0$ 对称条件的洛必达处理（3.3 节）；④ 完整写出问题 1 的数学问题（3.4 节）；⑤ 问题 4 的随体坐标 $\zeta=r/R(t)$ 变换与 $1/R^2$ 因子（9.1--9.7 节，涉及 Q4 时可推迟）。

\textbf{资料位置}：教程 2.1--3.4、9.1--9.7 节；传热学（第五版）第 2 章稳态热传导（p.51 起）圆筒壁导热部分——只读这一小段作对照。

\textbf{手推练习}：A1（平板热传导方程）、A9（$\zeta$ 变换路径 A 与路径 B）。

\textbf{学完标准}：能独立完成 A1；能说出三类边界条件的物理含义各举一例。

\subsection{模块 3\quad 解析解：分离变量与 Bessel 级数（P1，用到再学）}

\textbf{高中锚点}：三角函数叠加表示振动与波（物理）；幂级数（大一）。

\textbf{新学内容}：① 分离变量法全流程（教程 11.1 节，含特征方程 $\beta J_1(\beta)=\mathrm{Bi}\,J_0(\beta)$ 与系数 $A_n$ 的手推）；② 只需要"会用"的 Bessel 函数性质（$J_0'=-J_1$、正交权）；③ 级数求和与收敛判据。

\textbf{资料位置}：教程 11.1 节；传热学（第五版）第 3 章非稳态热传导（p.119 起）无限长圆柱级数解（作查阅，不作主线）；Crank《扩散的数学》。

\textbf{手推练习}：A11（Bessel 特征方程与 $A_n$）。

\textbf{学完标准}：能解释"比值 4"从哪来（A12）；知道解析解在论文里只用于验证而非求解。

\subsection{模块 4\quad 数值方法（P1，本项目的核心工具）}

\textbf{高中锚点}：数列递推（$a_{n+1}$ 由 $a_n$ 算出来）；速度 $\approx$ 位移差 $/$ 时间差。

\textbf{新学内容}：① 显式/隐式格式与稳定性手推（教程 4.4 节：$\xi=1-4r\sin^2(\theta/2)$、$\Delta t\le\Delta r^2/(2\alpha)$，Q1 上限 0.046 s）；② Crank--Nicolson 的梯形法则视角与截断误差 $-\Delta t^3u'''/12$（5.3--5.6 节）；③ 有限体积法：控制体体积、表面行、守恒恒等式（6.1--6.7 节）；④ 追赶法（7 章，模块 1 已铺垫）；⑤ 不动点/Picard 迭代与"为什么不用牛顿"（8 章）；⑥ 网格与步长的算账（12.2 节）。

\textbf{资料位置}：教程 4--8、12.2 节；传热学（第五版）第 4 章热传导问题的数值解法（p.171 起，差分格式与追赶法）；陶文铨《数值传热学》（进阶）。

\textbf{手推练习}：A3（显式稳定性）、A4（隐式放大因子）、A5（CN 截断误差）、A6（三个控制体体积）、A7（$N=5$ 的 CN 系统组装）、A12（收敛比值 4）。

\textbf{学完标准}：能独立推出 A3 并口头解释 Q1 为什么必须选隐式（0.046 s vs 实际 0.125 s）。

\section{五、学科二：物理（从高中热学往上接）}

\subsection{模块 5\quad 传导、扩散与对流边界（P0，必须先学）}

\textbf{高中锚点}：热传递三方式（初中：传导/对流/辐射）；扩散现象（高中选修 3-3：墨水入水）；比热容。

\textbf{新学内容}：① 傅里叶定律 $q=-k\partial T/\partial r$ 与菲克定律 $j=-D\partial C/\partial r$（教程 2.1 节，两条定律"同构"是本题全部数学的核心）；② 牛顿冷却公式型对流边界（3.2 节 Robin 条件的物理）；③ 数量级估计 $\sqrt{Dt}$（1.4 节，扩散深度 $\approx2.6$ mm）。

\textbf{资料位置}：教程 1.3、1.4、2.1、3.1--3.2 节；化工原理（第五版）上册 6.2.1 节傅里叶定律和热导率；传热学（第五版）第 5 章对流传热理论（p.207 起）只读牛顿冷却公式一段。

\textbf{手推练习}：A1（与模块 2 共用）。

\textbf{学完标准}：能不看书写出两条本构定律并说出每一项的单位；能解释"热方程与扩散方程长得一样"的物理原因。

\subsection{模块 6\quad 相变与湿空气（P1，读论文 7.6 前学）}

\textbf{高中锚点}：蒸发吸热、液化放热（初中/高中 3-3）；饱和汽压随温度升高而增大、相对湿度 $=$ 实际汽压/饱和汽压（3-3）；沸腾与蒸发的区别。

\textbf{新学内容}：① 湿空气参数：绝对湿度、相对湿度、焓（化工原理下册 14.2 节）；② 湿球温度：显热 $=$ 潜热的平衡式（化工原理下册 13.2.2 节式 (13-7)、(13-9)，$\alpha/k_H\approx1.09$ kJ/(kg$\cdot$$^\circ$C)）；③ 教程 10.4、10.5 节的湿球估计（约 40 $^\circ$C）与潜热数量级（$q_{\text{evap}}/q_{\text{sens}}\approx5$--$6$）。

\textbf{资料位置}：教程 10.4--10.5 节；化工原理（第五版）下册 13.2.2 节、14.2 节（湿空气状态参数）；上册第 7 章蒸发（7.3.3 节蒸发速率与传热温度差）。

\textbf{手推练习}：无独立练习；配套做 7.6 节的敏感性手推（A10 后半部分）。

\textbf{学完标准}：能写出湿球温度平衡式并解释"湿球温度是空气状态的函数，与物料无关"；能解释为什么 RH$\approx$57\% 对应湿球约 40--42 $^\circ$C。

\section{六、学科三：化学（从高中化学往上接）}

\subsection{模块 7\quad 含水率、活化能与扩散经验律（P1）}

\textbf{高中锚点}：质量分数与浓度（必修）；活化能定性概念：反应需要"翻越能垒"（选必 1）；温度升高反应加快。

\textbf{新学内容}：① 干基/湿基含水率的定义与换算（教程 10.1 节；论文符号 $C$ 为干基）；② Arrhenius 定量式与表观活化能手推（10.2 节：$E_a=3850\times8.3145\approx32.0$ kJ/mol）；③ 题给经验律 $D=D_0\mathrm{e}^{-0.45/C}$ 的化学含义：越干越难扩散（10.3 节）；④ 文献规律对照：实测 $D_{\mathrm{eff}}$ 随温度、含水率升高而增大（任迪峰论文 3.4 节结论，PDF 第 48--49 页）。

\textbf{资料位置}：教程 10.1--10.3 节；任迪峰论文（结论部分）。

\textbf{手推练习}：A10（$E_a=32.0$ kJ/mol）。

\textbf{学完标准}：给两组 $D$、$T$ 数据能反算出 $E_a$；能解释 $\mathrm{e}^{-0.45/C}$ 为什么表示"水分越少越难扩散"。

\subsection{模块 8\quad 平衡与敏感性（P2，了解即可）}

\textbf{高中锚点}：化学平衡（正逆速率相等时宏观不变）；勒夏特列原理。

\textbf{新学内容}：① 吸附/解吸平衡与平衡含水率概念（化工原理下册 14.2 节；论文 6.3 节"趋近平衡含水率"）；② 敏感性分析：单参数扰动 $\to$ 结果变化的定量估计（教程 10.6 节：$D$ 下降 32\% $\to$ 时长约 $+30\%$）；③ 有效参数观点（10.7 节，论文假设 H5 的逻辑链）。

\textbf{资料位置}：教程 10.6--10.7 节；化工原理（第五版）下册 14.2 节水分在气固两相间的平衡。

\textbf{手推练习}：A10 后半（敏感性因子 $0.68$）。

\textbf{学完标准}：能独立完成 10.6 节的指数敏感性手推并解释"为什么表面温度用 40 $^\circ$C 还是 50 $^\circ$C 只影响约 30\%"。

\section{七、学科四：工程科学（化工原理·固体干燥）}

\subsection{模块 9\quad 干燥两阶段与水分分类（P1，读论文 6.3 前学）}

\textbf{高中锚点}：湿衣服晾干先快后慢的生活经验；生物必修 1 的自由水/结合水；化学的饱和溶液。

\textbf{新学内容}：① 干燥速率曲线：恒速段（表面为非结合水覆盖、表面温度$=$湿球温度且不变）与降速段（四原因：实际汽化表面减小、汽化面内移、平衡蒸气压下降、内部扩散极慢），临界含水量（化工原理下册 14.3.1 节，图 14-12、表 14-1）；② 水分两套分类：\textbf{结合水分/非结合水分}（按结合方式，判据：蒸汽压是否等于 $p_s(T)$）与\textbf{平衡水分/自由水分}（按指定空气条件下能否除去）；关系：自由水分 $=$ 全部非结合水分 $+$ 部分结合水分，RH$=100\%$ 时平衡水分$=$结合水分；③ 文献实测结论：地黄薄层干燥\textbf{无恒速段}、干燥主要在降速段（任迪峰论文 p.101）；④ 对照论文 6.3 节：本文表面含水率 0.5--1.5 h 内持续下降、表面温度持续上升 $\Rightarrow$ 属"表面传质主导的快速干燥段"而非典型恒速段（术语规范结论）。

\textbf{资料位置}：教程 10.4、12.4 节；化工原理（第五版）下册 14.3.1 节；《External资料可用性评估报告》3.1、3.2 节；任迪峰论文结论章（p.101）。

\textbf{手推练习}：无独立练习；检查方式：能向同学讲清楚"自由水分 $\supset$ 非结合水分"并各举一例。

\textbf{学完标准}：能画出干燥速率曲线并标出预热/恒速/第一降速/第二降速四段与临界含水量；能说清干燥学的"自由水分"与生物学的"自由水"不是一回事。

\subsection{模块 10\quad 工艺原则与工程估算（P2，了解即可）}

\textbf{高中锚点}：晒干货怕高温变质的常识。

\textbf{新学内容}：① 工艺原则："恒速阶段表面温度维持在湿球温度，故即使易变质物料也允许高温；降速阶段物料温度逐渐升高，后期须防过热"（化工原理下册 14.3.1 节）——支撑论文 8.3 节变温推广；② 变温干燥机理：后期降温使湿度梯度与温度梯度同向，促进水分脱出、平均料温低有利保成分（任迪峰 p.101）；③ 间歇干燥时间工程估算式 $\tau_1=(G_c/A)(X_1-X_c)/N_A$（下册 14.3.2 节），与扩散模型的量级对照。

\textbf{资料位置}：化工原理（第五版）下册 14.3.1、14.3.2 节；任迪峰论文结论章；论文 8.3 节。

\textbf{学完标准}：能解释论文 8.3 节变温策略的每条依据出自哪条原理。

\section{八、编程（附录性模块，P2）}

\textbf{内容}：四个求解器的统一骨架——\texttt{thomas}（追赶法）、\texttt{crank\_step}（组装 CN 三对角系统）、\texttt{step}（交替 Picard 两轮）、主循环（推进 $+$ 按模板写 xlsx）与验证段（教程 12.3 节）。只需能读懂、会改参数重跑，不需要自己从零写。

\section{九、学习顺序与优先级总表}

\begin{center}
{\zihao{5}\begin{tabular}{p{3.4cm}p{1.1cm}p{4.2cm}p{2.6cm}p{3.4cm}}
\toprule
模块 & 优先级 & 对应教程节 & 对应论文位置 & 手推练习 \\
\midrule
1 微积分工具箱 & P0 & 1.1--1.5、7 章 & 5.4、附录 & A2、A8 \\
2 PDE 建模与定解条件 & P0 & 2--3 章、9 章 & 5.1--5.3 & A1、A9 \\
5 传导、扩散与对流边界 & P0 & 1.3--1.4、2.1、3.2 & 5.1、假设 H4 & A1 \\
4 数值方法 & P1 & 4--8 章、12.2 & 5.4、6.1--6.4 & A3--A7、A12 \\
3 解析解（Bessel） & P1 & 11.1 & 7.1 & A11 \\
7 含水率与活化能 & P1 & 10.1--10.3 & 5.1、7.5 & A10 \\
6 相变与湿空气 & P1 & 10.4--10.5 & 7.6 & A10 后半 \\
9 干燥两阶段与水分分类 & P1 & 10.4、12.4 & 6.3 & 讲给别人听 \\
8 平衡与敏感性 & P2 & 10.6--10.7 & 6.3、7.6、H5 & A10 后半 \\
10 工艺原则与工程估算 & P2 & 12.4 & 8.3 & 无 \\
编程 & P2 & 12.3 & 附录 & 无 \\
\bottomrule
\end{tabular}}
\end{center}

\textbf{建议节奏}：P0 三个模块按"5 $\to$ 1 $\to$ 2"的顺序学（先物理直觉、再数学语言、最后建模），约占总时间四成；之后按教程章序推进，每天 1--2 节教程 $+$ 当天对应练习，12 道练习全部手推过一遍即可进入论文逐段精读。

\section{十、资料索引与补充说明}

\textbf{底稿对应关系}：

\begin{enumerate}[leftmargin=2em]
\item \textbf{教程《A题方法详解》}：主线教材（12 章 $+$ 附录 A 练习 $+$ 附录 B 术语 33 条 $+$ 附录 C 学习路线）。
\item \textbf{论文《A题论文》}：目标文本，每章对应模块见上表。
\item \textbf{《External资料可用性评估报告》}：外部四份资料的对接点清单（第 3.2 节化工原理下册、3.1 节任迪峰论文与本图模块 6/9/10 直接对应）。
\item \textbf{外部教材只读节}：化工原理上册 6.2.1（傅里叶定律）、下册 13.2.2（湿球温度）、14.2--14.3（湿空气/干燥）；传热学（第五版）第 2--4 章对应小节（作查阅）；任迪峰论文结论章（p.101）。
\item \textbf{术语补充建议}：教程附录 B 建议新增 5 条——结合水分、非结合水分、平衡水分、自由水分、临界含水量（一句话定义见本图第三节速查表 3.3，展开见模块 9）。
\end{enumerate}

\textbf{最容易卡住的三处及桥接办法}：① 偏导与链式法则的符号恐惧——只看教程 1.1 节的"偏导 $=$ 固定其他变量求导"解释即可，不追求严格多元分析；② 分离变量法——教程 11.1 节已把每一步写全，先照抄一遍再合上书重推；③ Bessel 函数——只需记住 $J_0'=-J_1$ 与特征方程来历，不学函数论（教程 11.1 节已限定到这个范围）。

\textbf{说明}：传热学（第五版）为扫描版无文字层，只作翻阅对照；任迪峰论文数字多处乱码，只读其定性结论。所有外部资料只用于理解原理与验证，不改变论文的模型、参数与结果。

\end{document}
